% Part: lambda-calculus
% Chapter: syntax
% Section: eta

\documentclass[../../../include/open-logic-section]{subfiles}

\begin{document}

\olfileid[id]{lam}{syn}{eta}
\olsection{Konversi-$\eta$}

Terdapat relasi lain pada suku-$\lambda$. Dalam \olref[fv]{sec}, kita
menggunakan contoh $\lambd[x][(fx)]$, yang menerima suatu argumen dan
menerapkan $f$ padanya. Dengan kata lain, suku itu merupakan fungsi yang sama
dengan~$f$: $\lambd[x][(fx)]N$ dan $fN$ keduanya tereduksi menjadi~$fN$. Kita
menggunakan reduksi-$\eta$ (dan perluasan-$\eta$) untuk menangkap gagasan ini.

\begin{defn}[Kontraksi-$\eta$, $\eredone$]
  \ollabel{defn:beredone}
  \emph{Kontraksi-$\eta$} ($\eredone$) adalah relasi kompatibel terkecil pada
  suku-suku yang memenuhi syarat berikut:
  \[
    \lambd[x][M x] \eredone M \text{ asalkan } x \notin \FV{M}
  \]
\end{defn}

\begin{defn}[Reduksi-$\beta\eta$, $\bered$] \ollabel{defn:bered}
  \emph{Reduksi-$\beta\eta$} ($\bered$) adalah relasi refleksif dan transitif
  terkecil pada suku-suku yang memuat $\bredone$ dan $\eredone$, yakni aturan
  refleksivitas dan transitivitas ditambah dua aturan berikut:
  \begin{enumerate}
  \item Jika $M \bredone N$, maka $M \bered N$. \ollabel{defn:bered3}
  \item Jika $M \eredone N$, maka $M \bered N$. \ollabel{defn:bered4}
  \end{enumerate}

\end{defn}

\begin{defn}
  Kita memperluas relasi ekuivalensi $\equal$ dengan aturan konversi-$\eta$:
  \[
  \lambd[x][f x] \equal f
  \]
  dan menyatakan relasi yang diperluas tersebut sebagai $\equal[\eta]$.
\end{defn}

Ekuivalensi-$\eta$ penting karena berkaitan dengan ekstensionalitas suku
lambda:

\begin{defn}[Ekstensionalitas]
  Kita memperluas relasi ekuivalensi $\equal$ dengan aturan (\ext):
  \begin{center}
  Jika $Mx \equal Nx$, maka $M \equal N$, asalkan $x \notin \FV{MN}$.
  \end{center}
  Kita menyatakan relasi yang diperluas tersebut sebagai $\equal[\ext]$.
\end{defn}

Secara garis besar, aturan itu menyatakan bahwa dua suku, jika dipandang
sebagai fungsi, harus dianggap sama apabila keduanya berperilaku sama untuk
argumen yang sama.

Sekarang kita membuktikan bahwa aturan $\eta$ memberikan tepat
ekstensionalitas, dan tidak lebih dari itu.

\begin{thm}
  $M \equal[\ext] N$ jika dan hanya jika $M \equal[\eta] N$.
\end{thm}

\begin{proof}
  Pertama, kita membuktikan bahwa $\equal[\eta]$ tertutup di bawah aturan
  ekstensionalitas. Artinya, aturan \ext{} tidak menambahkan apa pun pada
  $\equal[\eta]$. Dengan demikian, $\equal[\eta]$ memuat $\equal[\ext]$, dan
  jika $M \equal[\ext] N$, maka $M \equal[\eta] N$.

  Untuk membuktikan bahwa $\equal[\eta]$ tertutup di bawah \ext, perhatikan
  bahwa bagi sebarang $M \equal N$ yang diperoleh dengan aturan \ext{}, kita
  mempunyai $Mx \equal[\eta] Nx$ sebagai premis. Maka kita memperoleh
  $\lambd[x][Mx] \equal[\eta] \lambd[x][Nx]$ berdasarkan salah satu aturan
  $\equal$; dengan menerapkan $\eta$ pada kedua ruas, kita memperoleh
  $M \equal[\eta] N$.

  Serupa dengan itu, kita membuktikan bahwa aturan $\eta$ termuat dalam
  $\equal[\ext]$. Bagi sebarang $\lambd[x][Mx]$ dan $M$ dengan
  $x \notin \FV{M}$, kita mempunyai
  $(\lambd[x][Mx])x \equal[\ext] Mx$, sehingga kita memperoleh
  $\lambd[x][Mx] \equal[\ext] M$ berdasarkan aturan \ext{}.
\end{proof}

\end{document}
